\section{Three translated factors and colored Tornheim jets}
\label{sec:triple}

The incoming reports ask whether three separated Stieltjes factors admit a
finite coordinate description in colored Tornheim or multiple Lerch jets.
We give an affirmative answer, with an explicit coefficient formula and
the local subtraction convention included.  The general relation between
Hurwitz products and Tornheim sums is classical
\cite{espinosa-tornheim,zhao-colored,zhao-zhou}; the contribution here is
the separated-shift, all-Stieltjes-index formulation and its regularization.

Write $e(x)=e^{2\pi i x}$.  For $|X|=|Y|=1$ with $X,Y\ne1$, define initially
in an absolutely convergent right half-space
\begin{equation}\label{eq:Tdef}
 \mathcal T(A,B,C;X,Y)
 =\sum_{m,n\ge1}\frac{X^mY^n}{m^A n^B(m+n)^C}.
\end{equation}
All powers of positive integers use their real logarithms.  A superscript
$\mathcal T^{(r,s,t)}$ below denotes derivatives in these three order
variables, with $X,Y$ fixed.

\begin{lemma}[Entire continuation for two nontrivial twists]
\label{lem:Tentire}
For the stated twists, $\mathcal T$ extends to an entire function of
$(A,B,C)\in\mathbb C^3$.  For $\Re C>0$ its continuation is
\begin{equation}\label{eq:Tmellin}
 \mathcal T(A,B,C;X,Y)=\frac1{\Gamma(C)}
 \int_0^\infty t^{C-1}\Li_A(Xe^{-t})\Li_B(Ye^{-t})\,dt.
\end{equation}
The integral is locally normally convergent in $A,B,C$ in that region.
At every integer $N\ge0$,
\begin{equation}\label{eq:Tnegative}
 \mathcal T(A,B,-N;X,Y)
 =\sum_{j=0}^{N}\binom Nj\Li_{A-j}(X)\Li_{B-N+j}(Y).
\end{equation}
\end{lemma}
\begin{proof}
In an initial region, insert the Gamma integral for $(m+n)^{-C}$ and
interchange absolutely convergent sums and the integral.  This gives
\eqref{eq:Tmellin}.  The polylogarithm is entire in its order at each fixed
argument on these two radial segments, and is holomorphic in $t$ near
zero because neither segment starts at $1$.  Consequently
\[
 f(t)=\Li_A(Xe^{-t})\Li_B(Ye^{-t})
      =\sum_{j=0}^{M}f_j(A,B)t^j+O_K(t^{M+1})
\]
uniformly on compact sets $K$ of orders.  The small neighborhood in $t$
can be chosen uniformly on $K$.  The expansion need not converge all the
way to $t=1$; its Taylor remainder still has the asserted local estimate,
which is all that subtraction requires.  Thus
\begin{align}\label{eq:Tcontinuation}
 \mathcal T(A,B,C;X,Y)=\frac1{\Gamma(C)}\bigg\{
 &\sum_{j=0}^{M}\frac{f_j(A,B)}{C+j}
 +\int_0^1 t^{C-1}\left(f(t)-\sum_{j=0}^{M}f_j(A,B)t^j\right)dt\notag\\
 &+\int_1^\infty t^{C-1}f(t)\,dt\bigg\}.
\end{align}
The subtracted integral is holomorphic for $\Re C>-M-1$.  The tail is
entire, since each polylogarithm decays exponentially, uniformly on
compact order sets.  The simple zeros of $1/\Gamma(C)$ cancel the apparent
poles at $C=0,-1,\ldots,-M$.  Increasing $M$ proves joint entire
continuation.  At $C=-N$, the value is $(-1)^N f^{(N)}(0)$.
Applying $d\Li_A(Xe^{-t})/dt=-\Li_{A-1}(Xe^{-t})$ and Leibniz's rule gives
\eqref{eq:Tnegative}.
\end{proof}

\subsection{A meromorphic identity before coefficient extraction}

Fix $a_1,a_2,a_3\in\mathbb R/\mathbb Z$, pairwise distinct.  Fractional
parts at isolated endpoints are immaterial in ordinary integrals.  Put
\begin{equation}\label{eq:Kdef}
 K(\boldsymbol\lambda;\boldsymbol a)
 =\int_0^1\prod_{j=1}^{3}
    \zeta(1-\lambda_j,\{x+a_j\})\,dx,
 \qquad \Re\lambda_j>0.
\end{equation}
For each $k\in\{1,2,3\}$, let $i<j$ be the other two indices and set
\[
 X_k=e(a_i-a_k),\qquad Y_k=e(a_j-a_k),\qquad
 \theta_k(\boldsymbol\lambda)=\frac\pi2(\lambda_i+\lambda_j-\lambda_k).
\]
The twists are nontrivial by the separation assumption.  Define
\begin{align}\label{eq:cones}
 S(\boldsymbol\lambda;\boldsymbol a)
 =\sum_{k=1}^{3}\bigl\{&e^{-i\theta_k(\boldsymbol\lambda)}
 \mathcal T(\lambda_i,\lambda_j,\lambda_k;X_k,Y_k)\notag\\
 &+e^{i\theta_k(\boldsymbol\lambda)}
 \mathcal T(\lambda_i,\lambda_j,\lambda_k;X_k^{-1},Y_k^{-1})\bigr\}.
\end{align}

\begin{theorem}[Six-cone identity]\label{thm:sixcone}
For $\Re\lambda_j>0$,
\begin{equation}\label{eq:sixcone}
 K(\boldsymbol\lambda;\boldsymbol a)
 =\frac{\prod_{j=1}^{3}\Gamma(\lambda_j)}
 {(2\pi)^{\lambda_1+\lambda_2+\lambda_3}}
 S(\boldsymbol\lambda;\boldsymbol a).
\end{equation}
The right side gives a meromorphic continuation to $\mathbb C^3$.
\end{theorem}
\begin{proof}
First assume $\Re\lambda_j>1$.  The nonzero Fourier coefficients of the
periodic Hurwitz function are
\begin{equation}\label{eq:hurwitz-fourier-lambda}
 \widehat{\zeta(1-\lambda,\cdot)}(n)
 =\frac{\Gamma(\lambda)}{(2\pi |n|)^\lambda}
    e^{-i\pi\lambda\operatorname{sgn}(n)/2},\qquad n\ne0,
\end{equation}
and the zero coefficient is zero \cite[Eq.~25.11.9]{dlmf-hurwitz}.
The three Fourier series converge absolutely, so integration keeps
exactly $n_1+n_2+n_3=0$ with every $n_j\ne0$.  Every such triple has either
two positive frequencies and one negative frequency, or the opposite
sign pattern.  In the first case the negative index is a unique $k$ and
the frequencies are $m,n,-m-n$.  Their phase and denominator give the
first term of the $k$th summand in \eqref{eq:cones}.  Reversing all signs
gives its second term.  This is a disjoint decomposition, so no factor
of two or ordering multiplicity is missing.

The integral in \eqref{eq:Kdef} is jointly holomorphic on
$\Re\lambda_j>0$.  Near a singular point of one factor, the other two
factors are smooth and uniformly bounded on compact order sets.  The
local majorants are $t^{\sigma-1}(1+|\log t|^M)$ with $\sigma>0$.
The right side is also holomorphic there by Lemma~\ref{lem:Tentire}.
The identity theorem extends the equality from the smaller right
half-space to the claimed domain.  The same formula supplies its
meromorphic continuation.
\end{proof}

\subsection{The precise endpoint finite part}

For $m\ge0$, use
\begin{equation}\label{eq:gamma-local}
 \gamma_m(t)=\frac{\log^m t}{t}+\gamma_m(1+t),\qquad 0<t<1.
\end{equation}
Let $x_k=\{-a_k\}$ and remove a right-hand arc of length $\varepsilon$
from each $x_k$, using the ordinary unit coordinate on the circle.
For sufficiently small $\varepsilon$ these arcs are disjoint.  Denote
the remaining set by $E_\varepsilon$, and put
\[
 A_k(\boldsymbol m;\boldsymbol a)
 =\prod_{j\ne k}\gamma_{m_j}(\{a_j-a_k\}).
\]

\begin{definition}[Separated triple finite part]\label{def:triplefp}
For $m_1,m_2,m_3\ge0$, define
\begin{align}\label{eq:triplefp}
 I_{\boldsymbol m}(\boldsymbol a)
 =\lim_{\varepsilon\downarrow0}\bigg\{
 &\int_{E_\varepsilon}\prod_{j=1}^{3}
      \gamma_{m_j}(\{x+a_j\})\,dx\notag\\
 &+\sum_{k=1}^{3}
 A_k(\boldsymbol m;\boldsymbol a)
       \frac{\log^{m_k+1}\varepsilon}{m_k+1}\bigg\}.
\end{align}
\end{definition}
The limit exists: after subtracting its displayed leading singular term,
the local product is $O(1+|\log t|^{m_k})$, hence integrable.  No product
of distributions with common singular support is used.

\begin{theorem}[All-index triple Stieltjes identity]
\label{thm:triple-jets}
Let
\begin{equation}\label{eq:Qtriple}
 Q(\boldsymbol\lambda;\boldsymbol a)
 =\prod_{j=1}^{3}\left\{\Gamma(1+\lambda_j)(2\pi)^{-\lambda_j}\right\}
       S(\boldsymbol\lambda;\boldsymbol a).
\end{equation}
This function is holomorphic near $\boldsymbol\lambda=0$, and
\begin{equation}\label{eq:triple-jets}
 \boxed{\quad
 I_{\boldsymbol m}(\boldsymbol a)
 =\left(\prod_{j=1}^{3}m_j!\right)
 [\lambda_1^{m_1+1}\lambda_2^{m_2+1}\lambda_3^{m_3+1}]
        Q(\boldsymbol\lambda;\boldsymbol a).
 \quad}
\end{equation}
Thus the answer uses finitely many derivatives of six colored Tornheim
kernels, of total order at most $m_1+m_2+m_3+3$.
\end{theorem}
\begin{proof}
In a local right coordinate $t$ at $x_k$, split the $k$th Hurwitz factor
as $t^{\lambda_k-1}+\zeta(1-\lambda_k,1+t)$.
For a smooth local multiplier $g$, analytic continuation of the singular
term is obtained from
\[
 \int_0^h t^{\lambda-1}g(t)\,dt
 =\frac{g(0)h^\lambda}{\lambda}
  +\int_0^h t^{\lambda-1}(g(t)-g(0))\,dt.
\]
The constant and higher Laurent coefficients of this formula are the
Hadamard finite parts of $\log^m(t)g(t)/t$, divided by $m!$.
Use disjoint neighborhoods of the three singularities and an ordinary
integral on their complement.  The other factors are smooth in the
local coordinate, although meromorphic in their own order variables.
Here is a uniform justification for extracting all three coefficients.
Let $g_k(t)$ be the product of those other factors, and define
\[
 F_\varepsilon(\boldsymbol\lambda)=
 \int_{E_\varepsilon}\prod_j\zeta(1-\lambda_j,\{x+a_j\})\,dx
 +\sum_k\frac{\varepsilon^{\lambda_k}g_k(0)}{\lambda_k}.
\]
On a small closed polydisk $|\lambda_j|\le\rho<1$, with $i,j\ne k$,
one has $\lambda_i\lambda_j(g_k(t)-g_k(0))=O(t)$ and
$\lambda_k\zeta(1-\lambda_k,1+t)=O(1)$ uniformly.  The local
decomposition therefore gives
\[
 \lambda_1\lambda_2\lambda_3(K-F_\varepsilon)
       =O(\varepsilon^{1-\rho})
\]
uniformly on that polydisk.  After the indicated pole factors are
removed, all these expressions are holomorphic.  Cauchy's coefficient
formula permits passage to the limit.  The counterterms contribute
exactly those in \eqref{eq:triplefp}, and hence
\[
 [\lambda_1^{m_1}\lambda_2^{m_2}\lambda_3^{m_3}]
 K=\frac{I_{\boldsymbol m}}{m_1!m_2!m_3!}.
\]
The convention $\zeta(1-\lambda,t)=-1/\lambda+
\sum_{m\ge0}\gamma_m(t)\lambda^m/m!$ explains why this equation has no
extra alternating signs.  Finally $K=Q/(\lambda_1\lambda_2\lambda_3)$
by $\Gamma(1+\lambda)=\lambda\Gamma(\lambda)$, proving
\eqref{eq:triple-jets}.
\end{proof}

\subsection{An explicit formula for three digamma factors}

Set $L=\gamma+\log(2\pi)$, where $\gamma=\gamma_0(1)$, and introduce
\[
 D_+=\partial-L-\frac{i\pi}{2},\qquad
 D_-=\partial-L+\frac{i\pi}{2}.
\]
The variable differentiated by each operator is specified by a subscript.
Since $\Gamma(1+\lambda)=1-\gamma\lambda+O(\lambda^2)$,
Theorem~\ref{thm:triple-jets} immediately gives
\begin{equation}\label{eq:triple-digamma}
 \boxed{\quad
 \FP\int_0^1\prod_{j=1}^{3}\psi(\{x+a_j\})\,dx
 =-2\Re\sum_{k=1}^{3}
 \left.
  D_{+,A}D_{+,B}D_{-,C}\mathcal T(A,B,C;X_k,Y_k)
 \right|_{A=B=C=0}.
 \quad}
\end{equation}
The finite part on the left is the negative of
$I_{(0,0,0)}$, because $\gamma_0=-\psi$.  This settles the specific
three-digamma coordinate question in the shifted-jets report.
It does not assert that these double-zeta derivatives reduce to ordinary
single zeta or Stieltjes values.

For example, for shifts $0,1/3,2/3$, write $\omega=e(1/3)$.  Symmetry of
$\mathcal T$ in $(A,X)$ and $(B,Y)$ shows that all three terms agree:
\begin{equation}\label{eq:triple-third}
 \FP\int_0^1\psi(x)\psi(\{x+\tfrac13\})
              \psi(\{x+\tfrac23\})\,dx
 =-6\Re\left.
 D_{+,A}D_{+,B}D_{-,C}\mathcal T(A,B,C;\omega,\omega^2)
 \right|_{A=B=C=0}.
\end{equation}
This is an exact root-of-unity specialization with a fully specified
regularization, not a numerical relation guess.

\subsection{A convergent representation of the first spectral derivative}

The first derivative in the final Tornheim order is especially convenient
for independent numerical checking.  With
$f(t)=\Li_A(Xe^{-t})\Li_B(Ye^{-t})$, equation
\eqref{eq:Tcontinuation} at $M=0$ gives
\begin{equation}\label{eq:Tfirst}
 \left.\partial_C\mathcal T(A,B,C;X,Y)\right|_{C=0}
 =\gamma f(0)+\int_0^1\frac{f(t)-f(0)}t\,dt
                    +\int_1^\infty\frac{f(t)}t\,dt.
\end{equation}
All $A,B$ derivatives pass through these convergent integrals on compact
sets.  Let
\[
 F_X(t)=\left.\partial_s\Li_s(Xe^{-t})\right|_{s=0}
            -(L+i\pi/2)\Li_0(Xe^{-t}),\qquad h(t)=F_X(t)F_Y(t).
\]
Then the quantity in one summand of \eqref{eq:triple-digamma} is
\begin{equation}\label{eq:triple-computable}
 \int_0^1\frac{h(t)-h(0)}t\,dt+
 \int_1^\infty\frac{h(t)}t\,dt+
       \left(-\log(2\pi)+\frac{i\pi}{2}\right)h(0).
\end{equation}
The Euler constant in the last coefficient has canceled against the
$\gamma f(0)$ term in \eqref{eq:Tfirst}.  Dropping that term would give
an incorrect finite constant.

\subsection{Integral inner orders and multiple polylogarithms}

For positive integers $A,B$ and initially sufficiently large $\Re C$,
partial fractions give the classical colored Tornheim reduction
\cite{zhao-colored}:
\begin{align}\label{eq:TtoMPL}
 \mathcal T(A,B,C;X,Y)
 ={}&\sum_{j=0}^{A-1}\binom{B+j-1}{j}
       \Li_{B+C+j,A-j}(Y,X/Y)\notag\\
 &+\sum_{j=0}^{B-1}\binom{A+j-1}{j}
       \Li_{A+C+j,B-j}(X,Y/X),
\end{align}
where $\Li_{p,q}(z,w)=\sum_{k>m\ge1}z^kw^m/(k^p m^q)$.
For completeness, multiply the elementary identity
\[
 \frac1{m^A n^B}
 =\sum_{j=0}^{A-1}\binom{B+j-1}{j}
       \frac1{m^{A-j}(m+n)^{B+j}}
 +\sum_{j=0}^{B-1}\binom{A+j-1}{j}
       \frac1{n^{B-j}(m+n)^{A+j}}
\]
by $X^mY^n(m+n)^{-C}$ and substitute $k=m+n$.
The displayed elementary identity follows inductively from
$1/(mn)=(1/m+1/n)/(m+n)$.
Analytic continuation then extends the complete expression in $C$.
At rational shifts, all the colors in \eqref{eq:TtoMPL} are roots of
unity.  One may differentiate this identity in $C$.  One must not
differentiate it in the discrete positive-integer parameters $A,B$;
it provides no such identity for those derivatives at zero.

\subsection{Ordinary triple log-Gamma correlations}

Let $G(x)=\log\Gamma(x)-\tfrac12\log(2\pi)=\zeta'(0,x)$.
Each translated $G$ has only a logarithmic singularity, so its triple
product is integrable.  Since $s_j=1-\lambda_j$, differentiating
\eqref{eq:Kdef} once in each $s_j$ at zero proves the ordinary identity
\begin{equation}\label{eq:triple-loggamma}
 \int_0^1\prod_{j=1}^{3}G(\{x+a_j\})\,dx
 =-\left.
 \partial_{\lambda_1}\partial_{\lambda_2}\partial_{\lambda_3}
 \left\{
  \frac{\prod_j\Gamma(\lambda_j)}{(2\pi)^{\sum_j\lambda_j}}
  S(\boldsymbol\lambda;\boldsymbol a)
 \right\}\right|_{\lambda_1=\lambda_2=\lambda_3=1}.
\end{equation}
Here no finite part is present.  More generally, any fixed spectral
derivatives at $s_j=-N_j\le0$ give convergent triple integrals of
Hurwitz jets, which include the normalized Stieltjes antiderivatives.
The local estimates are products of powers of $\log t$ times
nonnegative powers of $t$, so differentiation under the integral is
justified.  This extends the same six-coordinate mechanism from
singular Stieltjes functions to their ordinary integrated relatives.

Here is an explicit version at every primitive order.  For $r\ge1$ set
\begin{equation}\label{eq:Vprimitive}
 V_{m,r}(x)=(-1)^{m+1}m![u^{m+1}]
 \frac{\zeta(1-r+u,x)}{(r-1)!\prod_{j=1}^{r-1}(1-u/j)}.
\end{equation}
Empty products equal one.  These are finite linear combinations of
Hurwitz order derivatives at $1-r$ and have mean zero on the unit
interval.  In particular $V_{0,1}=-G$.
\begin{corollary}[All triple Stieltjes primitive correlations]
\label{cor:triple-primitives}
The pointwise identity $\partial_x^rV_{m,r}(x)=\gamma_m(x)$ holds on
$0<x<1$.  For pairwise distinct shifts and integers $r_j\ge1,m_j\ge0$,
\begin{align}\label{eq:triple-primitives}
 &\int_0^1\prod_{j=1}^{3}V_{m_j,r_j}(\{x+a_j\})\,dx\notag\\
 &\quad=(-1)^{m_1+m_2+m_3+3}\prod_{j=1}^{3}m_j!
 [u_1^{m_1+1}u_2^{m_2+1}u_3^{m_3+1}]
 \frac{K(\boldsymbol r-\boldsymbol u;\boldsymbol a)}
 {\prod_{j=1}^{3}\left\{(r_j-1)!
                   \prod_{h=1}^{r_j-1}(1-u_j/h)\right\}}.
\end{align}
In this formula $K$ is explicitly the six-cone right side of
\eqref{eq:sixcone}, and the integral is ordinary and absolutely
convergent.
\end{corollary}
\begin{proof}
The repeated Hurwitz derivative law gives
\[
 \partial_x^r\zeta(1-r+u,x)
 =-u(r-1)!\prod_{j=1}^{r-1}(1-u/j)\zeta(1+u,x).
\]
Extract the coefficient in \eqref{eq:Vprimitive} and use the defining
Stieltjes Laurent expansion.  This proves the primitive identity,
including its sign and normalization.  The identity
$\int_0^1\zeta(1-r+u,x)\,dx=0$ for small $u$ proves the asserted
mean-zero normalization.  Near $x=0$, write the zeta function as
$x^{r-1-u}+\zeta(1-r+u,1+x)$.  On $|u|\le\rho<1$, its derivatives
of any fixed order are bounded by a constant times
$1+x^{r-1-\rho}(1+|\log x|^M)$ for some finite $M$.
This is integrable even at $r=1$.  At each singularity of the product,
the other two shifted factors are smooth; the majorant is therefore
uniform on a sufficiently small polydisk of the $u_j$.
One can extract all coefficients under the integral, obtaining
\eqref{eq:triple-primitives}.
\end{proof}
